\subsection{作品目标}

针对上述安全问题与现有方案的不足，本作品旨在设计并实现一种面向外卖平台的隐私认证与可信审计系统——FoodShield。系统的核心安全目标如表\ref{tab:security-goals}所示。

\begin{table}[H]
\centering
\caption{FoodShield 系统安全目标}
\label{tab:security-goals}
\begin{tabular}{clp{8cm}}
\toprule
编号 & 安全目标 & 定义 \\
\midrule
G1 & 身份匿名性 & 通信参与方在日常业务过程中无法获取对方真实身份，仅能看到匿名标识或订单相关信息。 \\
G2 & 订单认证性 & 只有持有合法订单凭证的参与方才能进入对应订单的通信通道，攻击者无法伪造订单身份。 \\
G3 & 消息保密性 & 通信消息在存储阶段经过加密处理，管理员和未授权方默认无法查看消息明文。 \\
G4 & 日志完整性 & 通信日志的任何篡改、删除、插入或重排均可被检测，日志记录具备密码学完整性证明。 \\
G5 & 条件可溯源 & 在投诉或审计场景下，管理员可在权限控制和审计记录约束下基于订单号恢复必要身份信息。 \\
\bottomrule
\end{tabular}
\end{table}

上述五个安全目标相互支撑，共同构成了"日常匿名、异常可溯源"的可控匿名安全模型。其中，G1和G2保障日常通信中的隐私保护和认证安全，G3和G4保障通信数据的保密性和可信性，G5则在异常情况下为平台管理和纠纷处理提供追责能力。